\section{Comparison of Advantages and Shortcomings}
\label{sec:comparison}

\subsection{Cross-method comparison}
\label{subsec:cross-method-comparison}
Writing plan: this is the main synthesis section. Do not describe papers one by one. Compare assumptions, model inputs, supervision, output type, strengths, and weaknesses.

\begin{table}[t]
\centering
\small
\caption{Comparison of segmentation-centered method families.}
\label{tab:method-comparison-segmentation}
\begin{tabularx}{\linewidth}{p{0.20\linewidth} p{0.25\linewidth} X X}
\toprule
Method type & Representative papers & Advantage & Shortcoming \\
\midrule
CLIP-based OVSS & OVSeg, SAN, PnP-OVSS \citep{liang2023ovseg,xu2023san,luo2024pnpovss} & Strong open-vocabulary semantics and natural zero-shot category transfer. & CLIP is not designed for dense localization; spatial precision and prompt sensitivity remain issues. \\
Diffusion-based segmentation & ODISE \citep{xu2023odise} & Diffusion features provide rich spatial-semantic representations and support panoptic output. & Training and inference are costly, and the system is more complex than lightweight CLIP adapters. \\
SAM and VFM-based methods & SAM, ESC-Net, Trident \citep{kirillov2023sam,lee2025escnet,shi2025trident} & High-quality masks and strong generalization from promptable or compositional foundation models. & Native semantic understanding is limited; additional CLIP, VLM, or MLLM components are required. \\
\bottomrule
\end{tabularx}
\end{table}

\begin{table}[t]
\centering
\small
\caption{Comparison of generalist, reasoning, and extension-oriented method families.}
\label{tab:method-comparison-generalist}
\begin{tabularx}{\linewidth}{p{0.20\linewidth} p{0.25\linewidth} X X}
\toprule
Method type & Representative papers & Advantage & Shortcoming \\
\midrule
Unified decoder models & X-Decoder \citep{zou2023xdecoder} & Pixel, image, and language tasks can be handled in one framework. & Training is data-hungry and task balancing is difficult. \\
MLLM-based reasoning segmentation & LISA, GLaMM \citep{lai2024lisa,rasheed2024glamm} & Complex natural-language instructions, commonsense reasoning, and pixel grounding become possible. & Hallucination, fine boundary accuracy, heavy training cost, and evaluation ambiguity remain major problems. \\
Video, 3D, and complex-scene extensions & To be added & Closer to real applications with temporal or spatial consistency. & Annotation cost, memory, computation, and consistency evaluation are harder than in single-image settings. \\
\bottomrule
\end{tabularx}
\end{table}

\subsection{Trade-off axes}
\label{subsec:tradeoff-axes}
Writing plan: summarize the survey using several axes: semantic openness, mask quality, reasoning complexity, training cost, inference cost, and extensibility.

\subsubsection{Semantic openness versus spatial precision}
\label{subsubsec:semantic-vs-spatial}
Compare CLIP-based methods with SAM-based methods. CLIP supplies open semantics but weak localization, whereas SAM supplies masks but lacks category-level semantics \citep{liang2023ovseg,xu2023san,kirillov2023sam}.

\subsubsection{Training-based adaptation versus training-free composition}
\label{subsubsec:training-vs-training-free}
Compare OVSeg, SAN, and ESC-Net with PnP-OVSS and Trident. The former can learn task-specific structure, while the latter reduces training cost but may depend on heuristic fusion and multiple off-the-shelf models \citep{liang2023ovseg,xu2023san,lee2025escnet,luo2024pnpovss,shi2025trident}.

\subsubsection{Task-specific segmentation versus generalist reasoning}
\label{subsubsec:specialist-vs-generalist}
Compare specialized OVSS models with X-Decoder, LISA, and GLaMM. The main point is that generality improves interaction and reasoning ability but complicates training, evaluation, and deployment \citep{zou2023xdecoder,lai2024lisa,rasheed2024glamm}.

\placeholderfigure{Planned figure: capability trade-off map.}{fig:capability-tradeoff}{Figure plan: plot methods on axes such as semantic openness, mask precision, reasoning ability, and computational cost.}